\chapter{Traiter l'eau d'alimentation des chaudières}

L’eau qui alimente les chaudières doit être d’une qualité irréprochable afin d’éviter le phénomène de corrosion et l’encrassement à l’intérieur de la chaudière et des tuyauteries, vannes, pompes et autres équipements ou appareils de régulation. Les impuretés d’une eau, limpide d’apparence, peuvent être pour les chaudières et la tuyauterie la cause de graves inconvénients : dépôts de boue, tartre et corrosion.

\vspace{0.3cm}

L’objectif de ce module est de vous permettre de comprendre l’importance de la qualité de l’eau d’alimentation des chaudières en définissant toutes ses caractéristiques et les nombreux traitements indispensables à effectuer pour produire de la vapeur aux conditions souhaitées.



\begin{figure}[H]
\centering
\includegraphics[width=0.6\textwidth]{Images/M14/M14_1.png}
\end{figure}



\section{Composition de l’eau}


L’eau (H2O) est le résultat de l’union entre deux atomes d’hydrogène (H2) et un atome d’oxygène (O). Mais au cours de son cycle de vie, elle se charge d’impuretés qui peuvent être regroupés en trois catégories :

\begin{itemize}
\item  les matières en suspension (MES) ;
\item  les solides en solution ou sels dissous ;
\item  les gaz en solution ou gaz dissous.
\end{itemize}

\subsection{Les matières en suspension (MES)}

Il s’agit de boue, de sable, d’argile et de substances organiques (végétales ou animales). Elles ne doivent pas se retrouver dans le circuit de la chaudière, sinon elles provoquent des dépôts ou des phénomènes de moussage. 
Il convient donc de filtrer l’eau pour séparer les MES.


\subsection{Les sels dissous}

Ces sels proviennent de la dissolution des roches rencontrées par l’eau lors de son parcours.
Dans l’eau d’alimentation de chaudière il faut éviter la présence des éléments suivants :


\begin{itemize}
\item calcium ;
\item magnésium; 
\item sodium;
\item chlorures;
\item bicarbonates;
\item sulfates;
\item silicates formateurs de silice.
\end{itemize}


Ces éléments sont à l’origine de composés comme les tartres et la silice qui rendent l’eau « dure ». Dans les endroits les plus chauds, ils se déposent sur les parois des tuyauteries ou dans les ballons de chaudières. Il faut les éliminer à la source, par un procédé qui les extrait de l’eau alimentaire, ou empêcher leur agglomération en composé solide grâce aux produits chimiques de conditionnement.
\vspace{0.3cm}

Tous les sels dissous dans l’eau se présentent sous forme d’ions, qui sont des composés chimiques chargés positivement ou négativement.



\subsection{	Gaz dissous}

Les gaz dissous néfastes pour la chaudière sont le dioxyde de carbone (CO2) et l’oxygène (O2). Le premier est à l’origine du tartre, le second de la corrosion.
Il faut les éliminer par un dégazage et par des produits chimiques ajoutés à l’eau (conditionnement).


\section{	Qualité de l'eau de chaudière} 


La matière, sous toutes ses formes (gaz, liquide, solide), est constituée de molécules, qui sont des assemblages d’atomes.
L’atome peut s’apparenter à un nuage de petites billes – les électrons –, qui tournent autour d’une grosse boule – le noyau. L’image des planètes tournant autour du Soleil illustre bien cette idée.
Les électrons peuvent être comparés à des mains qui permettent aux atomes de se lier les uns aux autres. Le noyau est le corps de l’atome.
\vspace{0.3cm}

Ainsi les atomes, selon leur nombre de mains, peuvent s’associer avec plus ou moins de partenaires pour former des molécules.
\vspace{0.3cm}

Dans l’exemple de l’eau (H2O), l’oxygène possède deux mains et l’hydrogène n’en a qu’une. C’est pour cette raison que l’on trouve deux atomes d’hydrogène liés à un seul atome d’oxygène dans la molécule d’eau.
\vspace{0.3cm}


Les électrons ont une charge électrique négative (–), alors que les noyaux ont des charges positives (+).
Un atome n’est pas chargé électriquement car ils possèdent autant de charges négatives (–) – les électrons – que de charges positives (+) – le noyau.
\vspace{0.3cm}

Un ion ressemble à un atome qui a des électrons en moins ou en plus, donc il ressemble à un atome dont la charge est positive ou négative. Ils peuvent comme les atomes s’assembler pour former des molécules.
Les molécules d’eau peuvent se décomposer en deux composés, H+ et OH–, appartenant à la famille des ions : H+ et OH– s’assemblent pour donner H2O, et H2O se dissocie en H+ et OH–.
Dans l’eau presque tous les éléments sont des ions.
La qualité de l’eau dépend de la quantité de chacun des ions.


\subsection{	Concentration}

La concentration est le moyen de compter les substances chimiques, notamment dans une solution.
Son unité est une quantité (une masse ou un nombre) par unité de volume. Par exemple, 2 g de sel dans un litre d’eau représentent une concentration de sel dans l’eau de 2000 mg/l.
Nous utilisons le milligramme par litre (mg/l). Cette unité est extrêmement employée dans toutes les analyses de qualité d’eau.



\subsection{	Acides et bases}

Les acides, les bases et les sels constituent les trois principales catégories de composés chimiques qui interviennent à l’occasion du traitement de l’eau d’alimentation. En solution, ces composés se dissocient en particules chargées, qu’on appelle ions.
\vspace{0.3cm}

Les composés acides sont des composés qui contiennent de l’hydrogène, comme, par exemple l’acide chlorhydrique (HCl). Lorsqu’on met ce dernier en solution dans l’eau, l’acide se dissocie pour libérer des ions dérivés de l’hydrogène (l’ion hydronium, H+), qui fait qu’une solution est dite acide.
\vspace{0.3cm}

Les bases, appelées aussi alcalis, sont des composés qui contiennent des hydroxyles (OH–), comme par exemple l’hydroxyde de sodium (soude NaOH). Dans l’eau, ce dernier se dissocie pour libérer des ions hydroxyles (OH–). C’est cette libération d’ions hydroxyles qui fait qu’une solution est dite basique.
Les acides et les bases réagissent entre eux et se neutralisent. Il en résulte la formation d’un sel et de l’eau. Par exemple, le sel de cuisine est formé lorsque l’acide chlorhydrique (HCl) est neutralisée par une base, la soude (NaOH ou hydroxyde de sodium).
\vspace{0.3cm}

Ce sel en solution se dissocie en ions positifs (Na+) et en ions négatifs(Cl–). Les positifs s’appellent cations et les négatifs s’appellent anions.

\subsection{	pH} 

Le potentiel hydrogène (pH) indique le degré d’acidité d’une solution (voir également les modules n° 3 et n° 8). Il dépend de la concentration en H+ dans l’eau et il est mesuré sur une échelle allant de 0 à 14. 
Plus il y a de H+, plus le pH est bas (inférieur à 7), plus la solution est acide. 
Moins il y a de H+, plus le pH est grand (supérieur à 7), moins la solution est acide. 
S’il y a autant de H+ que de OH–, la solution est neutre, et le pH est égal à 7.

\vspace{0.3cm}

Remarque : Plus la température de l’eau s’élève, plus son acidité augmente (et donc plus elle est corrosive). En règle générale, le pH est donc mesuré à 20 °C.

\begin{figure}[H]
\centering
\includegraphics[width=0.6\textwidth]{Images/M14/M14_2.png}
\end{figure}

\subsection{Titres d’une eau}

Pour caractériser la qualité de l’eau de chaudière, nous utilisons des paramètres globaux appelés « titres ». Ils permettent de connaître la quantité des éléments chimiques importants présents dans l’eau, c’est-à-dire leur concentration.
\vspace{0.3cm}

Leur unité est le milliéquivalent par litre (meq/l) ou le degré français (°f). Cette unité représente la masse nécessaire d’un ion pour constituer une unité de charge électrique dans un litre d’eau.
Les trois principaux titres utilisés dans nos usines d’incinération sont :


\begin{itemize}
\item  le TA (titre alcalimétrique) ;
\item  le TAC (titre alcalimétrique complet) ;
\item  le TH (titre hydrotimétrique) ou dureté totale.

\end{itemize}

Le TA s’apparente à la somme des concentrations des carbonates (CO32–) et de l’hydroxyle (OH–).
Le TAC s’apparente à la somme des concentrations des carbonates (CO32–), de l’hydroxyle (OH–) et des bicarbonates (HCO3–).
Ces deux titres informent donc sur le pH et la possibilité de dépôts de tartre.
\vspace{0.3cm}

Alors que le TH s’apparente à la somme des concentrations du calcium (Ca2+) et du magnésium (Mg2+), responsables des dépôts de tartre.
Ils permettent de déterminer le pouvoir entartrant de l’eau et son acidité, les deux paramètres essentiels pour la qualité des eaux de chaudières.
Les valeurs guides (voir tableau index)  dépendent du type d’installation, mais il faut toujours avoir un TH proche de 0.


\subsection{Conductivité}

La conductivité représente la capacité de l’eau à laisser passer un courant électrique. Elle dépend directement de la quantité d’ions en solution car ils sont les uniques « transporteurs » du courant dans l’eau. Elle permet de mesurer la teneur de l’eau en sels dissous : plus une eau est salée, plus sa conductivité est grande.
\vspace{0.3cm}

La conductivité est pratique car elle renseigne très facilement et rapidement sur le pouvoir de corrosion de l’eau, ainsi que sur son pouvoir de déposer des solides.
Son unité est le micro Siemens par centimètre (µS/cm).
Plus l’eau est pure, plus sa conductivité est faible :


\begin{itemize}
\item eau pure : conductivité d'environ 0,1 µS/cm ;
\item eau hautement conductrice : conductivité supérieure à 500 µS/cm.
\end{itemize}


\section{	Phénomènes indésirables}

Un traitement de l’eau de nos chaudières est nécessaire pour lutter contre des phénomènes indésirables à l’origine d’usures prématurées et de détériorations importantes de notre équipement, mais aussi de perte de rendement.



\subsection{L’entartrage}

Sous l’effet de la température, certains sels, en particulier les sels de calcium et de magnésium, vont précipiter sous forme de dépôts durs et adhérents.
\vspace{0.3cm}

Ces dépôts vont se localiser préférentiellement sur les tubes les plus chauds. Ainsi ils diminuent les sections de passage pour les tubes des chaudières, provoquant une baisse d’irrigation des surfaces d’échange. Ils diminuent l’échange thermique et engendrent des surchauffes locales qui affaiblissent la résistance mécanique des aciers et peuvent provoquer à terme des déformations voire même leur rupture ou un coup de feu.
Par ailleurs, ces adhésions favorisent la corrosion en présence d’oxygène dissous.
\vspace{0.3cm}

Les boues non adhérentes se déposent en fond de chaudière pour les plus denses ou restent en suspension et peuvent être entraînées avec la vapeur (primage) et ainsi la polluer.

\subsection{Le primage}

C’est l’entraînement dans la vapeur de gouttes d’eau contenant des sels minéraux dissous ou des matières en suspension.
\vspace{0.3cm}

Le primage peut avoir pour origine :


 \begin{itemize}
\item  une mauvaise conception de la chaudière ;
\item  une variation brutale du débit de vapeur ;
\item  une mauvaise pression d’utilisation ;
\item  la présence d’une salinité excessive qui provoque un moussage ;
\item  la présence de MES ;
\item  la présence de matières organiques qui provoque aussi un moussage.


\end{itemize}


Le primage peut entraîner :


\begin{itemize}
\item  la pollution de la vapeur ;
\item  l’érosion et même la détérioration des tubes à cause de la circulation à grande vitesse de gouttelettes d’eau ;
\item  la baisse de rendement des échangeurs thermiques ;
\item  la baisse de la température de surchauffeur et risque de détérioration du surchauffeur ;
\item  l’encrassement de la turbine pouvant entraîner au mieux une perte de rendement, au pire le bris mécanique.


\end{itemize}

\subsection{	Corrosion} 

La corrosion correspond à un passage en solution du fer dans l’eau. Ceci se traduit par la formation de rouille pouvant conduire à la perforation de la tuyauterie.
Les surfaces d’échange des chaudières côté gaz de combustion et côté eau subissent des corrosions.
\vspace{0.3cm}

Pour le côté eau, les corrosions peuvent avoir pour origine :


\begin{itemize}
\item  un pH trop faible, inférieur à 7, donc une acidité trop forte qui provoque une « corrosion acide » ; c’est souvent le cas des canalisation de retour de condensats.
\item  un pH trop élevé, supérieur à 13, qui dissout la magnétite, oxyde naturellement protecteur de l’acier ; ce phénomène, rare et très localisé, est appelé « corrosion caustique ».
\item  la présence de gaz dissous tels que l’oxygène et le gaz carbonique qui proviennent de leur dissolution lors des contacts entre l’eau et l’air.

\end{itemize}


L’oxygène est l’un des principaux responsable des corrosions (par oxydation). Il y a apparition de piqûres très profondes de faible diamètre. Ces piqûres rongent le tube de l’intérieur.
L’oxygène est particulièrement néfaste pour les lignes de retour des condensats et pour toute la chaudière durant les arrêts si aucune mesure de conservation n’est prise.
\vspace{0.3cm}

Le dioxyde de carbone entraîné dans la vapeur provoque une acidification des condensats. Cela se traduit par un amincissement généralisé du métal.
\vspace{0.3cm}

Afin d’éviter ces phénomènes, il convient d’adapter la qualité de l’eau aux caractéristiques de la chaudière. Pour cela, il faut enlever les éléments indésirables (MES, sels dissous, gaz). C’est le rôle de l’épuration ou prétraitement. Il faut aussi réaliser un traitement en continu en ajoutant des réactifs appropriés capables de corriger chimiquement la qualité de l’eau.






\section{	Systèmes de traitement de l'eau}



Divers systèmes de traitement sont utilisés.
\begin{itemize}
\item  Pour limiter les dépôts, on cherche à éliminer le calcium, le magnésium et les minéraux présents dans l’eau. Le traitement consiste à utiliser des adoucisseurs d’eau ou des procédés de déminéralisation totale.
\item  Pour limiter la corrosion, on élimine l’oxygène de l’eau par dégazage ou/et par ajout d’additifs et on maintient l’eau d’alimentation à des pH basiques.
\item  Pour limiter le primage, on élimine les sels dissous en effectuant des purges de déconcentration de façon régulière.
\item  Pour rendre l’eau des retours de condensats non corrosive, on élimine le C02 par dégazage ou en ajoutant des réactifs dosés précisément.
\end{itemize}

Il existe trois principaux types de traitement de l'eau :

 \begin{itemize}
\item  l’épuration, qui consiste à éliminer les particules indésirables présentes dans l’eau ;
\item  le dégazage, qui consiste à éliminer les gaz présents dans l’eau ;
\item  le conditionnement, qui consiste à ajouter des réactifs chimiques à l’eau pour modifier ses caractéristiques.

\end{itemize} 

Les traitements sont variables d’un site à l’autre car soit la qualité de l’eau à traiter n’est pas la même (elle dépend en effet de la nature des sols), soit la vapeur n’a pas toujours la même qualité (dépend de son utilisation).


\subsection{	Purges} 

Les purges s’apparentent à une sorte d’épuration. Elles permettent de limiter le phénomène de concentration des sels.

Les pertes de vapeur, inévitables, peuvent provenir :


\begin{itemize}
\item  de l’utilisation même de la vapeur (injection de la vapeur dans le process) ;
\item  des purges d’eau réalisées à la pression atmosphérique avec donc une perte par revaporisation ;
\item  des fuites ;
\item  du dégazage.

\end{itemize}


Même si ces pertes de vapeur sont faibles, un appoint d’eau est nécessaire pour les compenser.

Quelle que soit la qualité de la déminéralisation réalisée sur l’eau d’appoint, celle-ci contient toujours des sels qui se concentrent en chaudière. De la même façon, lorsque vous faites bouillir de l’eau dans une casserole, il se forme toujours au fond et sur les parois des dépôts solides blanchâtres.
Pour ces raisons, un contrôle de la salinité est nécessaire.
\vspace{0.3cm}

On limite la salinité en réalisant une purge appelée purge de déconcentration (continue ou non, en fonction de la salinité ou de la consommation).
\vspace{0.3cm}

Lorsque les débits de purge nécessaires sont faibles, la purge peut être pratiquée de façon discontinue : un volume d’eau est évacué par manœuvre d’une vanne à ouverture rapide.
Le temps d’ouverture de la vanne est déterminé par :


 \begin{itemize}
\item  le volume d’eau que l’on désire évacuer pour maintenir une salinité basse en chaudière ;
\item  les intervalles de temps entre deux purges ;
\item  le diamètre de la vanne ;
\item  la pression de fonctionnement de la chaudière.

\end{itemize}


La purge discontinue peut être gênante :
\begin{itemize}
\item  du fait de l’appoint d’eau nécessaire pour la compenser, et de la chute de pression de la chaudière qui reçoit de l’eau à faible température ;
\item  du fait de la difficulté de gérer la récupération d’énergie sur les eaux de purge.

\end{itemize}

L’orifice de sortie est en général situé en fond de chaudière afin de réaliser dans le même temps la chasse des boues.
\vspace{0.3cm}

La purge continue peut être réalisée par asservissement d’une vanne à la résistivité (l’inverse de la conductivité) de l’eau en chaudière ou à débit constant et réglée manuellement si l’on connaît suffisamment bien le taux d’extraction.
\vspace{0.3cm}

Les purges ont pour conséquence :


\begin{itemize}
\item  de diminuer le rendement de production par la perte d’énergie contenue dans l’eau à la température de saturation de la vapeur. On réduit ces pertes par son utilisation pour le réchauffage de l’eau d’appoint entre autres ;
\item d’augmenter le volume d’eau d’appoint, donc la capacité de production nécessaire au traitement de cette eau.
\item d’augmenter la consommation de réactifs entraînés par les purges.

\end{itemize}


\begin{figure}[H]
\centering
\includegraphics[width=\textwidth]{Images/M14/M14_3.png}
\end{figure}

\vspace{0.3cm}
\subsection{Épuration}
\vspace {0,3cm}

\textbf{Élimination des sels dissous}
\vspace{0.3cm}

La technique la plus communément utilisée est un procédé par résines échangeuses d’ions. D’autres procédés, comme l’osmose inverse et la décarbonatation à la chaux, sont très rarement utilisés pour l’eau alimentaire des chaudières dans nos incinérateurs.

\vspace{1cm}
\textbf{Unité d'osmose inverse}
\vspace{0.3cm}

L’osmoseur est un filtre de très grande efficacité. Il laisse passer l’eau mais retient les particules et les sels dissous.
On obtient ainsi d’une part de l’eau très pure et d’autre part une eau chargée de toutes les impuretés, appelée aussi rétentat.


\begin{figure}[H]
\centering
\includegraphics[width=0.6\textwidth]{Images/M14/M14_4.png}
\end{figure}

\vspace{0.3cm}
\textbf{Le principe de l'osmose inverse}
\vspace {0.3cm}

Cette technique s’apparente à une filtration très poussée. Il s’agit de comprimer (la pression peut atteindre 70 bars) de l’eau pour qu’elle passe à travers une membrane « filtrante » appelée semi-perméable.
Dans le principe, l’installation est simple : elle comprend une membrane filtrante et une pompe haute pression.



\vspace{0.3cm}
\textbf{Equipement d'un osmoseur}
\vspace{0.3cm}

Dans la pratique, un osmoseur est constitué d’un module d’osmose, d’une pompe haute pression qui dirige l’eau déjà traitée sur la membrane du module. Avant de passer dans l’osmoseur, l’eau doit être déjà adoucie au minimum : si le module est alimentée en eau non adoucie, la membrane peut être détruite définitivement.
On a deux sorties : le rétentat ou rejet et le perméat (eau osmosée).
\vspace{0.3cm}

Des manomètres permettent la surveillance des pressions. Un pressostat arrête la pompe en cas de surpression ou en cas de manque d’eau.
La mise en œuvre de ce procédé pose différents problèmes : 

 \begin{itemize} 
\item il faut un traitement préalable de l’eau pour éviter le colmatage des membranes ;
\item les membranes se dégradent sous l’effet des fortes pressions, des micro-organismes bactériens ou de la chloration préventive.

\end{itemize}

\begin{figure}[H]
\centering
\includegraphics[width=0.6\textwidth]{Images/M14/M14_5.png}
\end{figure}


\subsection{Les résines échangeuses d’ions}

Ce sont des petites billes de plastique (Ø 0,4 à 1,2 mm) qui ont la capacité de retenir sélectivement certaines familles de sels dissous.
\vspace{0.3cm}

Les échanges ioniques
\vspace{0.3cm}

Le phénomène qui est à l’origine de la déminéralisation est l’échange ionique : on remplace les ions indésirables par d’autres acceptables du même signe (les anions par d’autres anions, les cations par d’autres cations). Cela n’est possible que grâce à l’existence de différences d’affinité entre les différents ions et la résine, c’est-à-dire que la résine préfère accrocher les ions indésirables en lieu et place de ceux qu’elle a initialement accroché à elle. Cette opération se déroule dans des échangeurs.

\vspace{0.3cm}
Les échangeurs
\vspace{0.3cm}

Un échangeur se présente comme un réservoir avec des entrées et des sorties. 
Le réservoir contient des billes de résine sur lesquelles viennent se fixer les composés indésirables au fur et à mesure que les composés acceptables sont relargués.
\vspace{0.3cm}

Après le passage d’une certaine quantité d’eau à traiter, la résine est saturée, c’est-à-dire qu’elle ne peut plus rien fixer. Alors elle subit un cycle de « nettoyage » ou plutôt de régénération, avant de redevenir disponible pour le cycle de traitement.

\vspace{0.3cm}
Les cycles nécessitent donc une régulation.
\vspace{0.3cm}

Il existe deux types de résines :

\begin{itemize}
\item  les résines échangeuses de cations retiennent exclusivement les sels dissous porteurs de charges (+), donc le calcium (Ca2+) et le magnésium (Mg2+) qui sont à l’origine du tartre.
\item les résines échangeuses d’anions retiennent exclusivement les sels dissous porteur de charges (–), donc les silicates (HSiO3–), les carbonates(CO32–), et les bicarbonates (HCO3–), à l’origine aussi du tartre et de la silice.
\end{itemize}


Elles se nomment échangeuses car pour retenir les sels dissous indésirable elles « relarguent » dans l’eau un composé chimique qui ne pose pas de problèmes pour nos chaudières : il y a un échange d’ions.


En fonction de la qualité d’eau de chaudière attendue, il est possible de choisir parmi trois procédés d’épuration utilisant les résines :
\begin{itemize}
\item  l’adoucissement ;
\item  la décarbonatation puis adoucissement ;
\item  la déminéralisation totale.

\end{itemize}





\textbf{Adoucissement}
\vspace{0.3cm}

L’adoucissement correspond à un simple retrait des calcium et magnésium de l’eau par des résines échangeuses de cations qui relarguent du sel de sodium.
\vspace{0.3cm}

L’adoucisseur est composé :

\begin{itemize}
\item  d’une cuve contenant les résines échangeuses de sels ;
\item  d’un bloc de distribution ou panoplie de vannes ;
\item d’un bac à sel contenant une réserve d’eau fortement salée appelée saumure.

\end{itemize}


\begin{figure}[H]
\centering
\includegraphics[width=0.6\textwidth]{Images/M14/M14_6.png}
\end{figure}


\begin{figure}[H]
\centering
\includegraphics [width=0.6\textwidth]{Images/M14/M14_7.png}
\end{figure}

Le principe de fonctionnement se décompose en deux phases :

\begin{itemize}
\item  l’adoucissement ;
\item  la régénération (injection de sodium sur la résine saturée en calcium et magnésium).

\end{itemize}

La régénération est décomposée en quatre phases :
\begin{itemize}
\item  le détassage ;
\item  le rechargement des résines en sels ou saumurage ;
\item  le rinçage lent ;
\item  le rinçage rapide.
\end{itemize}

Dans la phase d’adoucissement, l’eau brute est envoyée sur des résines qui retiennent le calcium, l’eau est renvoyée vers l’installation avec le sodium.
\vspace{0.3cm}

Les résines une fois saturées en calcium ne sont plus efficaces. Il faut alors les régénérer.
Le passage de l’eau de haut en bas finit par tasser les résines, il faut alors les détasser en envoyant de l’eau de bas en haut. Cette opération permet ainsi un brassage des résines et leur homogénéisation.


\begin{figure}[H]
\centering
\includegraphics[width=\textwidth]{Images/M14/M14_8.png}
\end{figure}


Pour remplacer le calcium fixé par les résines, on injecte dans l’adoucisseur une solution de saumure (sel de cuisine) qui passe au travers de l’adoucisseur. Le calcium est remplacé par du sodium.
\vspace{0.3cm}

Le rinçage permet au sodium de réagir complètement.
\vspace{0.3cm}

Pour chasser la solution de saumure qui ne s’est pas fixée sur la résine, cette dernière est rincée par de l’eau brute. Pendant cette opération, une solution de saumure est à nouveau préparée pour recommencer un nouveau cycle.
\vspace{0.3cm}

La décision de la régénération peut être prise :

\begin{itemize}
\item  en fonction d’une durée de fonctionnement en service ;
\item  en fonction du volume écoulé ;
\item  en optimisant les deux données précédentes.

\end{itemize}


\vspace{0.3cm}
\textbf{Décarbonatation-adoucissement} 
\vspace{0.3cm}


La décarbonatation se déroule dans une tour, dans laquelle on pulvérise l’eau pour libérer le CO2 qui y est dissout.
Ainsi on a minimisé les possibilités de fabriquer des carbonates donc du tartre.


\begin{figure}[H]
\centering
\includegraphics[width=0.5\textwidth]{Images/M14/M14_9.png}
\end{figure}


\begin{figure}[H]
\centering
\includegraphics[width=0.5\textwidth]{Images/M14/M14_10.png}
\end{figure}

\vspace{0.3cm}
\textbf{Déminéralisation totale}
\vspace{0.3cm} 

La déminéralisation permet une élimination de tous les sels dissous. Ce procédé s’est généralisé sur toutes les usines d’incinération récentes, à cause des performances demandées à la chaudière.
La déminéralisation de l’eau consiste à enlever tous les composés susceptibles de constituer des minéraux ou des sels insolubles.
\vspace{0.3cm}

Il s’agit d’un adoucissement plus poussé qui fonctionne sur le même principe.
Cette eau est très agressive vis-à-vis des métaux, car elle est tellement pure qu’elle a tendance à dissoudre le métal ; les canalisations doivent être en inox ou en matière plastique.
\vspace{0.3cm}

On alimente la station de déminéralisation par l’eau à traiter.
Cette eau va passer à travers quatre éléments principaux dans la chronologie suivante :


\begin{itemize}
\item  un échangeur de cations ;
\item  une tour de décarbonatation ;
\item  un échangeur d’anions ;
\item  un échangeur à lit mélangé ou échangeur mixte dans quelques cas.

\end{itemize}

a) Retrait des cations dans un échangeur de cations
\vspace{0.3cm}

Les cations ou ions [+] sont fixés sur de la résine chargée [–] car les opposés s’attirent.
Les ions indésirables arrêtés sont Ca2+, Mg2+.
Le calcium et le magnésium sont retirés de l’eau. On a donc « enlevé » la dureté. C’est pour cette raison que cette phase se nomme aussi adoucissement.
\vspace{0.3cm}

En contrepartie, ces ions ont été remplacés par des ions H+ acceptables.
Avec cette résine, on a capté tous les ions [+] indésirables.
Donc, dans l’eau, après passage dans la résine, il ne reste plus que des ions [–] à traiter dans une autre résine.


b) Décarbonatation
\vspace{0.3cm}

Il s’agit du même principe que pour l’adoucissement : la concentration des carbonates et bicarbonates est diminuée.


c) Retrait des anions dans un échangeur d’anions
\vspace{0.3cm}

L’eau décarbonatée, contenant essentiellement des ions H+ comme ions [+] et différents ions [–] dont les chlorures (Cl–), les sulfates (SO42–), les bicarbonates (HCO3–) et les silicates (HSiO3–), passe dans un second type de résine chargée [+].
Donc tous les ions chargés [–] sont retenus par la résine.
Cette étape permet de diminuer la concentration des acides les plus « corrosifs » de l’eau.
De même, elle permet de diminuer la quantité de silice. 

En contrepartie, ces ions ont été remplacés par des ions OH– (acceptables).
À ce stade, on voit que l’eau ne contient plus que H+ et OH– ( en théorie).


d) Retrait des derniers ions [+] et [–] par polishing dans certains cas
\vspace{0.3cm}

Cette étape n’existe pas dans toutes les usines.
L’eau qui a déjà été bien purifiée subit un dernier passage dans un échangeur mixte.
Il y a un mélange des deux types de résine, une qui capte les éléments [–] et une qui capte les éléments [+].
Ainsi, on obtient un haut niveau de pureté .


e) Stockage
\vspace{0.3cm}

L’eau purifiée est acheminée dans une bâche de stockage.


Le fonctionnement
\vspace{0.3cm}


a) Les durées
Le fonctionnement de l’installation est cyclique. Il y a une période de « service » pendant laquelle l’eau est déminéralisée.
Et il y a aussi une période de « régénération » pendant laquelle la résine est reconditionnée pour une nouvelle utilisation.


b) Principe de la déminéralisation



\begin{figure}[H]
\centering
\includegraphics[width=0.8\textwidth]{Images/M14/M14_11.png}
\end{figure}

\vspace{0.3cm} 


Les phases de service et de régénération
\vspace{0.3cm}


a) Les échangeurs de cations et d’anions
\vspace{0.3cm}

Quatre étapes se succèdent lors du fonctionnement normal de ce type d’échangeur : une étape de service, qui permet la production d’eau, et trois petites étapes de reconditionnement.
1. Étape de service.
2. Étape d’introduction d’un excès de régénérant (acide ou base selon l’échangeur).
3. Étape de déplacement pour répartir le régénérant.
4. Étape de rinçage pour enlever l’excès de régénérant qui ne s’est pas lié à la résine.






\begin{figure}[H]
\centering
\includegraphics[width=0.8\textwidth]{Images/M14/M14_12.png}
\end{figure}


 
b) L’échangeur à lit mélangé
\vspace{0.3cm}

Son principe est le même que celui des échangeurs précédents.


\begin{figure}[H]
\centering
\includegraphics[width=\textwidth]{Images/M14/M14_13.png}
\end{figure}





\newpage

Légende du schéma ci-contre : 

\begin{enumerate}
\item L’eau à traiter entre dans l’échangeur à résine cationique, qui fixe exclusivement le calcium, le magnésium et le sodium. Ils sont fixés parce qu’ils ont une charge électrique positive.
\item  Ces ions sont remplacés dans l’eau par l’ion H+. L’eau qui ressort de l’échangeur cationique ne contient plus le calcium, le magnésium et le sodium et va dans la tour de décarbonatation (9).
\item  Cette eau en entrant dans la tour de décarbonatation perd le dioxyde de carbone qui y était dissous. Ce gaz est évacué par un évent (10). Et l’eau qui ne contient plus de dioxyde de carbone dissous est aspirée par la pompe (11) pour être envoyée dans l’échangeur anionique.
\item  Dans l’échangeur anionique, les ions négatifs tels que le bicarbonate, le sulfate, le silicate et le chlorure sont fixés sur la résine. Donc à la sortie de l’échangeur l’eau ne les contient plus.
\item Ces ions négatifs sont remplacés dans l’eau par les ions OH–.
\item  Les ions H+ qui viennent de l’échangeur cationique réagissent (7) avec les ions OH– pour former de l’eau, H2O.
\item  Quand les résines cationiques ont fixé le calcium, le magnésium et le sodium, il faut les régénérer. De l’acide est injecté dans l’échangeur cationique.
\item  De même pour l’échangeur anioniques, qui avait capté le bicarbonate, le sulfate, le silicate et le chlorure, dans lesquelles on injecte de la soude.
\item  et 14 Lors de l’injection des régénérants, les résines relarguent tous les ions indésirables, le bicarbonate, le sulfate, le silicate et le chlorure pour les ions négatifs et le calcium, le magnésium et le sodium pour les ions positifs. Ces ions s’associent pour former des sels qu’on rejette dans les eaux usées après une neutralisation du pH.

\end{enumerate}

\subsection{Élimination des gaz dissous}

\vspace{0.3cm}
\textbf{Principe}
\vspace{0.3cm}

Plus la température est grande moins les gaz se dissolvent facilement dans l’eau(plus leur solubilité est petite), jusqu’à une température de 110 °C.
Par exemple environ 9mg d’oxygène se dissolvent dans un litre d’eau à 20 °C.
Alors qu’il n’y en a plus que 0,18 mg (50 fois moins) à 99 °C.
Le principe du dégazage thermique consiste donc à porter l’eau à une température d’environ 105 °C (température optimale) en faisant une injection de vapeur dans l’eau de la bâche dégazante.
L’eau ainsi dégazée, dont la teneur résiduelle en oxygène est proche de 0,03mg/l est stockée dans la bâche alimentaire.
Au-delà de 110 °C, la solubilité de l’oxygène augmente à nouveau, il est donc inutile de travailler à plus haute température.




\begin{figure}[H]
\centering
\includegraphics[width=\textwidth]{Images/M14/M14_14.png}
\end{figure}




\section{Conditionnement}


Le conditionnement de l’eau de chaudière consiste à ajouter des produits chimiques pour éviter les risques d’entartrage, de corrosion, d’embouage des circuits d’eau et même de primage. Il vient en complément des traitements d’épuration que nous venons de voir.
D’ailleurs, les produits chimiques de conditionnement seront d’autant plus efficaces que l’épuration sera bien faite.


\subsection{Les différents produits chimiques}

\begin{itemize}
\item  Produits dispersants qui évitent l’agglomération des particules en suspension dans l’eau.
\item  Produits sequestrants qui ralentissent la formation de solides dans l’eau.
\item  Produits complexants qui piègent les composés du tartre pour empêcher sa formation,
\item  Produits alcalins qui augmentent le pH évitant, ainsi la corrosion acide et le dépôt de silice.
\item  Produits réducteurs d’oxygène qui détruisent l’oxygène dissous, ainsi la corrosion est fortement limitée.
\item  Inhibiteurs de corrosion par action filmogène qui fabriquent une peau fine sur les canalisations d’eau, limitant ainsi fortement la corrosion.

\end{itemize}

Les sociétés qui vendent ces produits de conditionnement proposent des mélanges savamment dosés.


Les produits essentiels au conditionnement de l’eau de chaudières sont au nombre de trois :

\begin{itemize}
\item  les phosphates pour disperser les boues de calcaire. Ils sont injectés dans les ballons de chaudières ;
\item  les alcalinisants qui évitent la corrosion acide dans le retour des condensats. Ils appartiennent classiquement à la famille des amines, molécules qui contiennent de l’azote. Ils sont souvent injectés avec les phosphates ;
\item  les réducteurs d’oxygène qui détruisent l’oxygène gazeux, évitant ainsi la corrosion par oxydation. Ils sont injectés au niveau de la bâche dégazante.
\end{itemize}



\subsection{La surveillance de la qualité des eaux de chaudière}

Pour établir la nature et la quantité d’impuretés que contient l’eau de la chaudière ou l’eau d’alimentation, il est nécessaire de prélever des échantillons de cette eau et de les soumettre à divers essais ou analyses.
Les flacons à échantillons doivent être propres et rincés à l’eau à échantillonner. Avant de prélever un échantillon, il faut laisser couler cette eau un certain temps afin d’éviter d’utiliser l’eau qui aurait séjourné trop longtemps dans le conduit. On doit laisser refroidir cette eau à la température de la pièce et filtrer les matières en suspension ou les laisser décanter.
On doit procéder aux analyses aussitôt que possible après le prélèvement de l’échantillon.

Sur certaines installations, il y a des échantillonneurs automatiques qui permettent l’analyse en continu.


La surveillance de la qualité de l’eau se fait en plusieurs points :

\begin{itemize}
\item en continu sur l’eau d’appoint en sortie de chaîne de déminéralisation On mesure la résistivité ou la conductivité de l’eau déminéralisée. La régénération est lancée si la résistivité est trop faible ;
\item  à l’aide de prise d’échantillons sur l’eau alimentaire en sortie du dégazeur. On mesure la teneur en oxygène. L’injection du réducteur d’oxygène sera augmentée dans le cas d’une concentration en oxygène élevée. La teneur résiduelle en réducteur (aquarex 1340) doit être comprise entre 0,05 et 0,5 mg/l ;
\item  à l’aide de prise d’échantillon sur l’eau de ballon de chaudière en sortie de purge. Si la concentration en sel est trop importante, on augmente le débit de purge. Si le pH est trop faible, on augmente le débit d’injection des produits alcanisants. Le dosage optimum en phosphatant (aquarex 1210) est de l’ordre de 300 g/m3 en chaudière.
\item  à l’aide de prise d’échantillon d’eau sur les retours de condensats pour s’assurer que le pH n’est pas trop bas (donc acide). Sinon il faut le remonter en augmentant la concentration des amines volatiles. Le dosage moyen d’amines volatiles (aquarex 1120) est de l’ordre de 2g/°f TAC.m3 d’eau alimentaire

\end{itemize}

\vspace{0.3cm}
\textbf{Les différentes techniques d'analyse}
\vspace{0.3cm}

La détermination du TH, du TA, du TAC et des chlorures se fait le plus souvent par méthode titrimétrique qui se déroule ainsi :


\begin{itemize}
\item ajouter une petite quantité de réactif coloré à l'échantillon ;
\item ajouter des gouttes d'un autre produit chimique, tout en remuant, jusqu'à ce que la couleur change.
\end{itemize}

La quantité du produit chimique ajouté pour le changement de couleur révèle le taux de l’impureté recherchée.

\vspace{0.3cm}

La détermination du taux, de réducteur d’oxygène, des phosphatants et de la silice restant en chaudière est le plus souvent réalisée par une méthode photo-colorimétrique. Cette technique consiste à mesurer les concentrations en exposant les échantillons au passage d’un rayon de couleur choisie en fonction du produit à mesurer. Ce rayonnement est absorbé uniquement par les éléments recherchés. L’intensité lumineuse du rayonnement sera d’autant plus basse que la concentration en éléments recherchés dans l’eau sera grande.
La modification de la luminosité est observée par un capteur lumineux qui permet d’afficher directement la concentration des éléments recherchés.
\vspace{0.3cm}

Le pH et la conductivité sont déterminés par des appareils électroniques munis de deux électrodes réunies dans un petit cylindre en verre ou en plastique. Ce petit cylindre est plongé dans la solution à tester et l’affichage de l’appareil donne la valeur du paramètre recherché.




\begin{figure}[H]
\centering
\includegraphics[width=0.4\textwidth]{Images/M14/M14_15.png}
\end{figure}

\section{Sécurité}


Certains procédés de traitement d’eau nécessitent la manipulation de certains réactifs dangereux pour l’organisme si l’on ne prend aucune précaution.
Les acides attaquent la peau et provoquent de graves brûlures, lorsqu’il y a un contact direct. Les bases ont des effets proches des acides mais plus graves.
Pour la manipulation de ces produits, il faut absolument utiliser des gants en matière plastique, des lunettes, un tablier et des bottes pour se protéger des projections.
\vspace{0.3cm}

Des douches sont disponibles à proximité des postes de traitement. En cas de projection, il faut les utiliser en lavant à grande eau la partie touchée.
\vspace{0.3cm}

Les acides peuvent s’évaporer sous l’effet de la chaleur ou s’ils sont mélangés avec d’autres produits.
Les produits sous forme de poudre sont volatils et peuvent s’introduire dans les voies respiratoires. Ils peuvent être dangereux pour les voies respiratoires.
L’emploi d’un masque s’avère parfois nécessaire.
\vspace{0.3cm}

Dans tous les cas, l’étiquetage sur les produits signale quelles précautions à prendre pour leurs manipulations et les risques encourus.


\begin{figure}[H]
\centering
\includegraphics[width=0.5\textwidth]{Images/M14/M14_16.png}
\end{figure}


Interprétation des contrôles et consignes d’exploitation


\begin{figure}[H]
\centering
\includegraphics[width=\textwidth]{Images/M14/M14_17.png}
\end{figure}


Tableau récapitulatif des valeurs guides pour la qualité de l'eau


1. Eau déminéralisée 

Eau d'alimentation
Pression : 40-60 bar
pH : > 8,5
TH (°f) : < 0,05
02 : < 0,02 mg/l


Eau de chaudière
TAC (°f) : < 15
TA (°f) : > 0,5 TAC (essentiel)
pH : 10 à 11
phosphates (mg/l) : 10 à 60
Max. silice (mg/l) : 15




2. Eau adoucie décarbonatée ou adoucie 

Eau d'alimentation
Pression : 35-40 bar
pH : > 8,5
TH (°f) : < 0,10
02 : < 0,02 mg/l


Eau de chaudière
Max. TAC (°f): 40
Min. TAC (°f) : 10
pH : 11,2
phosphates (mg/l) : 30
Max. silice (mg/l): 40


